\begin{lemma}[The $x$-derivative of the Lebesgue kernel]
Fix $s > 0$ and write $r = e^{-s} \in (0,1)$ and $\sigma = 1 - e^{-2s} = 1-r^2 \in (0,1)$,
and let
\[
  p_s(x,y) = \frac{1}{\pi^{n/2}\sigma^{n/2}}
    \exp\Bigl(-\frac{|rx-y|^2}{\sigma}\Bigr)
\]
be the Lebesgue kernel of \noderef{MehlerLebesgueKernel}.  Then for every
$x, y \in \mathbb{R}^n$ and every $j \in \{1,\ldots,n\}$ the function
$w \mapsto p_s(w,y)$ is differentiable at $x$, and
\[
  D_j^x\,p_s(x,y)
    = \sqrt2\,r\,\frac{y_j - r x_j}{\sigma}\;p_s(x,y) ,
\]
where $D_j = 2^{-1/2}\partial_j$ is the energy-normalized partial derivative, so
that the left-hand side means $2^{-1/2}\,\partial_{x_j}p_s(x,y)$.

\emph{The normalization matters and is the easiest thing here to misread.}  In
the Euclidean normalization the same computation reads
$\partial_{x_j}p_s = 2r(y_j-rx_j)\sigma^{-1}p_s$, which is the form the reference
paper uses when it derives the normalized kernel; the $\sqrt2$ above is that $2$
multiplied by the $2^{-1/2}$ of $D = 2^{-1/2}\nabla$, the convention of
\noderef{GaussD} that the whole tablet uses.  The two displays are the same
identity, not two different constants.

The statement is also exactly the appendix's derivative identity in different
letters.  With the appendix's reparametrization $\sigma = 1-r^2$, $d = (1-r)x$,
$z = y-x$ and $q_\sigma(z) = (\pi\sigma)^{-n/2}e^{-|z|^2/\sigma}$ one has
$z + d = y - rx$, hence $p_s(x,y) = q_\sigma(z+d)$ and
$\sqrt2\,r\,(y_j-rx_j)\sigma^{-1}p_s(x,y) = \sqrt2\,r\,V_{j,\sigma}(z+d)$ with
$V_{j,\sigma}(z) = (z_j/\sigma)q_\sigma(z)$.  So the appendix's displayed
derivative identity is this node, and the appendix's $(\sigma, d, z, q_\sigma)$
notation is a change of letters to be introduced locally where it is used, not a
second fact.

Unlike \noderef{MehlerLebesgueKernel}, this is an assertion at \emph{every}
$(x,y)$: no almost-everywhere qualifier appears, because $p_s$ is an explicit
smooth function of its first argument and no equivalence class is involved.
\end{lemma}
\begin{proof}
Fix $s>0$, $y \in \mathbb{R}^n$ and $j$.  Since $s>0$ we have $0 < r < 1$, hence
$0 < \sigma = 1-r^2 < 1$; in particular $\sigma \neq 0$, so every division by
$\sigma$ below is legitimate, and the constant
$c = \pi^{-n/2}\sigma^{-n/2}$ is a well-defined positive real.

\emph{Step 1: the exponent is a polynomial.}
Write $A(w) = -\sigma^{-1}|rw-y|^2$, so that $p_s(w,y) = c\,e^{A(w)}$.  In
coordinates,
\[
  A(w) = -\sigma^{-1}\sum_{i=1}^n (r\,w_i - y_i)^2 ,
\]
a polynomial in the coordinates $w_1,\ldots,w_n$; polynomials are differentiable
everywhere, and the $j$th partial derivative is computed termwise.  Only the
$i = j$ term depends on $w_j$, and by the chain rule for $t \mapsto (rt-y_j)^2$,
\[
  \partial_{w_j}A(w)
    = -\sigma^{-1}\cdot 2\,(r w_j - y_j)\cdot r
    = \frac{2r\,(y_j - r w_j)}{\sigma} ,
\]
the last step using $-(rw_j - y_j) = y_j - r w_j$.

\emph{Step 2: differentiating the exponential.}
The real exponential is differentiable with derivative itself, so by the chain
rule $w \mapsto c\,e^{A(w)}$ is differentiable at every $w$, in particular at
$x$, which is the first assertion.  Its $j$th partial derivative is
\[
  \partial_{w_j}\bigl(c\,e^{A(w)}\bigr)
    = c\,e^{A(w)}\,\partial_{w_j}A(w)
    = p_s(w,y)\cdot\frac{2r\,(y_j - r w_j)}{\sigma} .
\]
Evaluating at $w = x$,
\begin{equation*}
  \partial_{x_j}p_s(x,y)
    = \frac{2r\,(y_j - r x_j)}{\sigma}\,p_s(x,y) . \tag{1}
\end{equation*}
This is the Euclidean-normalized form quoted in the statement, and it is the
identity the reference paper displays when it passes from the semigroup to the
normalized kernel.

\emph{Step 3: passing to the $D$ normalization.}
By \noderef{GaussD} the energy-normalized derivative is $D = 2^{-1/2}\nabla$, so
$D_j^x p_s(x,y) = 2^{-1/2}\partial_{x_j}p_s(x,y)$.  Multiplying $(1)$ by
$2^{-1/2}$ and using $2^{-1/2}\cdot 2 = 2^{1/2} = \sqrt2$,
\[
  D_j^x p_s(x,y)
    = 2^{-1/2}\cdot\frac{2r\,(y_j-rx_j)}{\sigma}\,p_s(x,y)
    = \sqrt2\,r\,\frac{y_j - r x_j}{\sigma}\,p_s(x,y) ,
\]
which is the assertion.
\end{proof}
